% ============================================================
%   CHAPTER 3: Architecture Concepts
% ============================================================

\chapter{Architecture Concepts}
\label{chap:architecture_concepts}

\begin{table}[H]
\caption{Revision History – Chapter 3}
\centering
\begin{tabular}{llll}
\toprule
Version & Date & Author & Description \\
\midrule
1.0 & \today & Y. Sangale & Initial Draft \\
\bottomrule
\end{tabular}
\end{table}

This section defines the fundamental architectural principles of the 64-bit \gls{riscv} single-core microcontroller. It provides a high-level technical description of the \gls{cpu} core, memory subsystem, bus structure, interrupt design, clocking strategy, and register architecture. 

\textcolor{blue}{\textbf{\textit{Concepts related to power management, security, and debug will be fully defined in the final specification revision.}}}

\subsection{Technology}

The microcontroller is conceptually targeted for implementation in a modern \textbf{90\,nm CMOS technology node}. 
This node provides a balanced compromise between integration density, power efficiency, and manufacturing cost, 
making it suitable for compact embedded microcontrollers and low-complexity \gls{SoC} designs.

The 90\,nm node enables:
\begin{itemize}
  \item Higher operating frequencies compared to older planar CMOS nodes.
  \item Reduced dynamic and static power consumption.
  \item Dense \gls{sram} bitcells allowing larger memories in a small silicon footprint.
\end{itemize}

Although the conceptual design and functional model are developed and validated using synthesizable \gls{hdl} on \gls{fpga} platforms, 
the microarchitecture, clocking structures, and memory system were designed with 90\,nm \gls{asic} fabrication constraints in mind.
This includes fully synchronous logic, controlled critical paths, and modular blocks suitable for automated place-and-route flows.

\subsection{\gls{cpu} Concept}
The \gls{cpu} core is a compact, single instruction \gls{rv64i} processor optimized for determinism and simplicity.

\begin{itemize}
    \item \textbf{Pipeline:} A simple scalar \gls{pipeline} (Fetch–Decode–Execute–Writeback) without speculation or out-of-order features.
    \item \textbf{\gls{isa} Support:} Fully compliant with \gls{rv64i} + M-extension (integer multiply/divide).
    \item \textbf{Datapath Width:} 64-bit \gls{alu}, registers, and internal buses.
    \item \textbf{Register File:} 32 × 64-bit general-purpose registers (x0–x31), with direct \gls{cpu} access (no bus).
    \item \textbf{Memory Access Model:} \gls{Harvard} architecture — separate \gls{I-MEM} and \gls{D-MEM} buses.
          \begin{itemize}
              \item Dedicated instruction interface to \gls{I-MEM}.
              \item Dedicated data interface to \gls{D-MEM}.
              \item Memory-mapped I/O via Wishbone for all peripherals.
          \end{itemize}
    \item \textbf{Program Counter:} 64-bit \gls{pc} with generation logic (\gls{pc} +4, branch target, jump target).
    \item \textbf{Interrupt Integration:} Receives a single aggregated interrupt line from the Interrupt Controller 
          and vectoring is performed as per \gls{riscv} rules.
\end{itemize}

\begin{itemize}
    \item \textbf{Pipeline:} Single-stage or simplified multi-stage \gls{pipeline} (Fetch → Decode → Execute → Writeback).
    \item \textbf{Data Width:} 64-bit \gls{alu}, register file, and datapath.
    \item \textbf{Instruction Execution:} One instruction executed at a time; no out-of-order execution.
    \item \textbf{\gls{isa}Support:} \gls{rv64i} with full support for:
        \begin{itemize}
            \item 64-bit integer arithmetic and logical operations,
            \item Load/store with byte, half-word, word, and double-word support,
            \item Multiply/divide (M-extension),
            \item Branch and jump instructions.
        \end{itemize}
    \item \textbf{Interrupt Input:} Single external interrupt line fed by the Interrupt Controller.
    \item \textbf{Memory Access:} \gls{Harvard} architecture — separate \gls{I-MEM} and \gls{D-MEM} buses.
\end{itemize}

\subsection{Bus Concept}
The microcontroller employs a simple and predictable communication architecture centered around a 
\textbf{single-master bus topology}. The \gls{cpu} core is the only bus master and is responsible for initiating 
all data and peripheral access transactions. This simplifies arbitration, reduces latency, and provides 
deterministic behavior suitable for education, prototyping, and low-complexity embedded applications.

The system uses two categories of interconnects:

\begin{itemize}
    \item \textbf{Direct \gls{cpu} Connections}  
    The following modules are connected directly to the \gls{cpu} datapath and do not use the system bus:
    \begin{itemize}
        \item \textbf{Instruction Memory (\gls{I-MEM}):} Accessed through a dedicated instruction fetch interface.
        \item \textbf{Data Memory (\gls{D-MEM}):} Connected through a dedicated data access path for load/store instructions.
        \item \textbf{Register File (RF):} Integrated directly into the \gls{cpu} \gls{pipeline}.
    \end{itemize}

    These direct connections enable single-cycle access to critical datapath components and support 
    a clean \gls{Harvard} architectural split.

    \item \textbf{Wishbone-Based Peripheral Bus}  
    All other memory-mapped peripherals are connected using a \textbf{Wishbone-compliant} slave interface.  
    This bus is used for:
    \begin{itemize}
        \item General Purpose I/O (\gls{gpio}),
        \item Universal Asynchronous Receiver/Transmitter (\gls{uart}),
        \item Quad-SPI Master (\gls{qspi}),
        \item Interrupt Controller,
    \end{itemize}

    The Wishbone bus provides standardized signaling for address, data, strobes, and acknowledgments, enabling:
    \begin{itemize}
        \item Modular and plug-and-play peripheral integration,
        \item Predictable, single-master timing behavior,
        \item Straightforward extension with new peripherals,
        \item Clean separation between \gls{cpu} datapath and peripheral logic.
    \end{itemize}
\end{itemize}


\subsection{Memory Concept}
The memory subsystem is based on a clean \gls{Harvard} architecture with separated instruction and data paths.

\begin{itemize}
    \item \textbf{Instruction Memory (\gls{I-MEM}):}
          \begin{itemize}
              \item 64\,KB read-only memory for program storage.
              \item Dedicated \gls{cpu} instruction fetch path (non-Wishbone).
              \item High-performance single-cycle access for most operations.
          \end{itemize}

    \item \textbf{Data Memory (\gls{D-MEM}):}
          \begin{itemize}
              \item 64\,KB data RAM for stack, heap, and program data.
              \item Direct \gls{cpu} data access (non-Wishbone).
              \item Fully synchronous with \gls{cpu} core.
          \end{itemize}

    \item \textbf{Memory-Mapped I/O:}
          \begin{itemize}
              \item Peripherals mapped in the 0x1000\_0000 address region.
              \item Accessed exclusively via the Wishbone peripheral bus.
          \end{itemize}

    \item \textbf{Extensibility:}
          \begin{itemize}
              \item Optional I-Cache and D-Cache may be added in future without changing the CPU–bus interface.
          \end{itemize}
\end{itemize}

\subsection{Interrupt Concept}
The interrupt subsystem is based on a configurable multi-source interrupt controller supplying a single interrupt input to the \gls{cpu}.

\begin{itemize}
    \item \textbf{\gls{IRQ} Lines:} 8 independent inputs (\texttt{\gls{IRQ}[7:0]}) with per-line configuration.
    \item \textbf{Trigger Types:}
        \begin{itemize}
            \item Level-high
            \item Rising edge
            \item Falling edge
            \item Both edges
            \item Pulse/one-shot
        \end{itemize}
    \item \textbf{Edge Detection:} Implemented using per-line previous-state registers.
    \item \textbf{Control Registers:} 
          Mask, type-select, and pending-status registers accessible via Wishbone.
    \item \textbf{Vector Computation:}
          \[
              \text{ISR Vector} = \text{BASE\_ADDR} + (\text{\gls{IRQ}\_ID} \times 0x10)
          \]
    \item \textbf{Priority Model:} Fixed priority (IRQ0 highest) for the initial revision.
    \item \textbf{\gls{cpu} Interface:} Outputs a single interrupt line to the \gls{cpu}'s machine-mode interrupt system.
\end{itemize}

\subsection{Clock Concept}
The microcontroller uses an 80\,MHz main system clock with optional peripheral clock division.

\begin{itemize}
    \item \textbf{Core Clock:} 80\,MHz clock feeding \gls{cpu}, \gls{I-MEM}, and \gls{D-MEM}.
    \item \textbf{Peripheral Clocks:} Generated via a programmable divider block:
          \begin{itemize}
              \item Divide-by-2, /4, /8, etc.
              \item Per-peripheral selectable clock domains.
          \end{itemize}
    \item \textbf{Clock Control Registers:} Accessible via Wishbone.
    \item \textbf{Clock Domain Crossing:} 
          Standard two-flip-flop synchronizers used where required (e.g., \gls{uart}, \gls{qspi}).
\end{itemize}

\subsection{System Control Concept}
System control is centered around a simple and robust global reset architecture.  
The microcontroller implements an \textbf{external \gls{gpio}-based reset input}, which provides a direct and intuitive 
mechanism to reinitialize the entire system.

\begin{itemize}
    \item \textbf{External Reset Input:}  
          A dedicated \gls{gpio} pin is configured as a reset input. When driven \textbf{low}, it asserts a 
          system-wide synchronous reset.

    \item \textbf{Reset Behavior:}
          \begin{itemize}
              \item All \gls{cpu} state (\gls{pc} , registers, CSRs) is reset to defined initial v\gls{alu}es.
              \item Instruction Memory and Data Memory interfaces are reset.
              \item All Wishbone peripherals (\gls{gpio}, \gls{uart}, \gls{qspi}, \gls{IRQ} controller, Clock Control, JTAG region) 
                    return to default configuration.
              \item All internal state machines return to IDLE.
          \end{itemize}

    \item \textbf{Synchronization:}  
          The external reset signal passes through a two-stage synchronizer to ensure safe integration into the 
          main clock domain.

    \item \textbf{Reset Release:}  
          When the external \gls{gpio} reset input returns high, the system exits reset synchronously at the next 
          active clock edge.

\end{itemize}

This minimal reset architecture provides deterministic system initialization and simplifies external control 
during prototyping and debugging.

\subsection{Register Concept}
The Register Concept defines the structure and behavior of all \textbf{memory-mapped peripheral registers} 
accessible through the Wishbone bus. These registers are used to configure, control, and monitor the operation 
of \gls{gpio}, \gls{uart}, \gls{qspi}, the Interrupt Controller, the Clock Control Unit, and other on-chip peripherals.

\begin{itemize}
    \item \textbf{Scope:}  
          Only peripheral registers are included.  
          \gls{cpu}-internal structures such as the Register File and CSRs are part of the \gls{cpu} microarchitecture 
          and are not memory-mapped.

    \item \textbf{Register Types:}
        \begin{itemize}
            \item Control registers (enable, mode select, configuration)
            \item Status registers (interrupt flags, error states, peripheral activity)
            \item Data registers (\gls{uart} TX/RX, \gls{qspi} data buffer, \gls{gpio} input/output)
            \item Configuration registers (\gls{IRQ} trigger type, \gls{gpio} direction, clock divisor)
        \end{itemize}

    \item \textbf{Access Method:}  
          All registers are accessed via standard 32-bit or 64-bit \gls{cpu} load/store instructions through the 
          Wishbone data bus.

    \item \textbf{Address Alignment:}  
          All registers are aligned to 32-bit word boundaries.

    \item \textbf{Reset Behavior:}  
          On system reset, all registers revert to defined default states to ensure predictable system startup.

\end{itemize}


\subsection{Power Concept}

\textcolor{blue}{\textbf{\textit{Section to be completed in the final document revision.}}}

\subsection{System Protection and Security Concept}

\textcolor{blue}{\textbf{\textit{Section to be completed in the final document revision.}}}

\subsection{Debug Concept}

\textcolor{blue}{\textbf{\textit{Section to be completed in the final document revision.}}}


